% !TeX root = ../机器学习教材（李航）笔记.tex
% 由主文档通过 \input 载入；请编译主文档。
\setcounter{chapter}{1} % 保留原书章号。
\chapter{感知机}
\label{chap:lh-perceptron}

给定带有两类标签的训练样本，最直接的分类思路是寻找一个超平面，使两类样本分别位于它的两侧。\emph{感知机（perceptron）}采用这一线性判别规则，并通过逐次纠正误分样本来确定超平面。它提供了理解监督学习的一条完整路线：先规定预测模型，再设计可计算的损失，最后研究更新算法在什么条件下能够终止。这里的收敛性是对有限训练集的分离保证，不直接给出未知数据上的分类误差保证。

全章使用训练集 $T=\{(x_i,y_i)\}_{i=1}^{N}$，其中 $N\geq 1$，$x_i\in\mathbb R^d$ 为第 $i$ 个样本的特征向量，$y_i\in\{-1,+1\}$ 为其类别标签。向量均视为列向量，$u^\top v$ 表示欧氏内积，$\|u\|$ 表示欧氏范数。样本下标写在右下角，向量坐标写为 $x^{(j)}$，以免混淆样本编号与特征编号。

\section{感知机模型}
\label{sec:lh-perceptron-model}

\begin{definition}[感知机模型]
\label{def:lh-perceptron}
设输入空间 $\mathcal X\subseteq\mathbb R^d$，输出空间 $\mathcal Y=\{-1,+1\}$。给定权重向量（weight vector）$w\in\mathbb R^d$ 和偏置（bias）$b\in\mathbb R$，感知机的得分函数和分类函数分别为
\begin{equation}
\label{eq:lh-perceptron-model}
s_{w,b}(x)=w^\top x+b,
\qquad
f_{w,b}(x)=\operatorname{sign}\bigl(s_{w,b}(x)\bigr).
\end{equation}
本笔记将符号函数在零点的取值约定为
\[
\operatorname{sign}(z)=
\begin{cases}
+1,&z\geq 0,\\
-1,&z<0.
\end{cases}
\]
这一约定只决定恰落在决策边界上的输入如何预测；训练时另行规定边界样本仍须触发更新。
\end{definition}

感知机直接由特征给出类别，属于\emph{判别模型（discriminative model）}。其假设空间是
\[
\mathcal H=
\left\{x\longmapsto\operatorname{sign}(w^\top x+b):
w\in\mathbb R^d,\ b\in\mathbb R\right\}.
\]
其中 $s_{w,b}$ 是实值得分，$f_{w,b}$ 才是取两个类别值的分类器；二者不能混作同一个函数。得分本身没有被建模为概率，也不能直接解释为分类置信概率。

当 $w\neq 0$ 时，方程
\begin{equation}
\label{eq:lh-perceptron-hyperplane}
S=\{x\in\mathbb R^d:w^\top x+b=0\}
\end{equation}
定义一个超平面（hyperplane），其法向量为 $w$，两个开半空间分别满足 $s_{w,b}(x)>0$ 和 $s_{w,b}(x)<0$。正方向由 $w$ 指向得分增加的一侧。对任意 $c>0$，$(w,b)$ 与 $(cw,cb)$ 给出同一分类器，因此参数通常不唯一。$w=0$ 时得分为常数，不能将式~\eqref{eq:lh-perceptron-hyperplane} 解释为通常的超平面，但这样的参数可作为算法的初始状态。

训练阶段利用标签确定 $w,b$；预测阶段则固定所得参数，对新输入 $x$ 计算式~\eqref{eq:lh-perceptron-model}。训练集能否被分开，首先取决于它相对于当前特征表示是否线性可分，而不是取决于训练算法运行了多少次。

\section{感知机学习策略}
\label{sec:lh-perceptron-strategy}

\subsection{数据集的线性可分性}
\label{subsec:lh-perceptron-separability}

\begin{definition}[线性可分性]
\label{def:lh-linear-separability}
若存在 $w\in\mathbb R^d$、$b\in\mathbb R$，使训练集中每个样本都满足
\begin{equation}
\label{eq:lh-perceptron-separable}
y_i(w^\top x_i+b)>0,
\qquad i=1,\ldots,N,
\end{equation}
则称 $T$ 为\emph{线性可分数据集（linearly separable data set）}；否则称为线性不可分。满足上述不等式的超平面称为分离超平面（separating hyperplane）。
\end{definition}

式~\eqref{eq:lh-perceptron-separable} 将两类样本的条件合为一个式子：$y_i=+1$ 时要求得分严格为正，$y_i=-1$ 时要求得分严格为负。这里强调严格不等式，因而没有训练样本落在分离超平面上。当两类样本均存在时，分离向量必有 $w\neq0$；若训练集只有一类，则可以用常数偏置完成分类。

\begin{definition}[函数间隔与几何间隔]
\label{def:lh-perceptron-margin}
对参数 $(w,b)$，第 $i$ 个样本的\emph{函数间隔（functional margin）}为
\begin{equation}
\label{eq:lh-perceptron-functional-margin}
m_i(w,b)=y_i(w^\top x_i+b).
\end{equation}
当 $w\neq0$ 时，其有符号几何间隔（geometric margin）为 $m_i(w,b)/\|w\|$，点 $x_i$ 到超平面的无符号距离为 $|w^\top x_i+b|/\|w\|$。
\end{definition}

上述距离公式可由正交分解得到：从 $x_i$ 沿法向量方向移动至超平面，垂足为
\[
x_i-\frac{w^\top x_i+b}{\|w\|^2}w,
\]
移动长度就是 $|w^\top x_i+b|/\|w\|$。函数间隔为正表示样本处于正确一侧，为负表示误分，为零表示落在边界上。正比例缩放参数会改变函数间隔，但不改变几何间隔。

\subsection{从误分类到感知机损失}
\label{subsec:lh-perceptron-loss}

直接最小化误分类数需要优化不连续的类别判断，不利于利用梯度信息。教材从误分点到分离超平面的距离引出替代目标：对于 $m_i(w,b)<0$ 的点，其到超平面的距离为 $-m_i(w,b)/\|w\|$。去掉公共归一化因子 $1/\|w\|$ 后，可得到只惩罚负函数间隔的感知机损失。\emph[caution]{由于 $\|w\|$ 也随参数改变，去掉该因子是选择新的损失函数，并非保持原距离最小化问题不变的代数化简。}

\begin{definition}[感知机损失]
\label{def:lh-perceptron-loss}
单个样本的感知机损失和训练集的总损失分别定义为
\begin{equation}
\label{eq:lh-perceptron-loss}
\ell_i(w,b)=\max\{0,-y_i(w^\top x_i+b)\},
\qquad
L(w,b)=\sum_{i=1}^{N}\ell_i(w,b).
\end{equation}
若将误分或位于边界的样本指标集记为
\[
M(w,b)=\{i:y_i(w^\top x_i+b)\leq0\},
\]
则等价地有
\begin{equation}
\label{eq:lh-perceptron-loss-active}
L(w,b)=-\sum_{i\in M(w,b)}y_i(w^\top x_i+b).
\end{equation}
边界样本在和式中的贡献为零。采用平均损失 $L/N$ 只改变整体尺度，不改变其最小值点集合。
\end{definition}

若 $w\neq0$，全部误分点到超平面的距离之和为 $L(w,b)/\|w\|$，而 $L$ 本身是负函数间隔之和。不同参数之间不能只凭误分点个数判断 $L$ 的大小：误分点较少时，某个点的负间隔仍可能很大。反过来，参数整体缩小时，$L$ 可以减小而分类结果完全不变。

\begin{proposition}[损失的基本性质]
\label{prop:lh-perceptron-loss-properties}
函数 $L$ 非负、连续、凸且分段线性，但一般不处处可导。它满足
\begin{equation}
\label{eq:lh-perceptron-zero-loss}
L(w,b)=0
\quad\Longleftrightarrow\quad
y_i(w^\top x_i+b)\geq0\quad\text{对所有 }i.
\end{equation}
严格分离蕴含零损失，零损失一般不蕴含严格分离。
\end{proposition}

\begin{proof}
每个 $\ell_i$ 是常数函数 $0$ 与仿射函数 $-y_i(w^\top x_i+b)$ 的最大值，因此非负、连续、凸且分段线性，有限求和保留这些性质。两个分支交界处的斜率一般不同，所以不能据此声称函数处处可导。有限个非负项之和为零当且仅当每项均为零，这就得到式~\eqref{eq:lh-perceptron-zero-loss}。严格正间隔当然满足非负间隔；反例是 $(w,b)=(0,0)$，此时任何训练集都满足 $L(0,0)=0$，但所有间隔均为零。
\end{proof}

教材将学习策略写为
\begin{equation}
\label{eq:lh-perceptron-objective}
\min_{w\in\mathbb R^d,\ b\in\mathbb R} L(w,b).
\end{equation}
理解这一表达时必须同时保留算法的更新和停止规则：由于零参数总是式~\eqref{eq:lh-perceptron-objective} 的最小值点，单独求解这个无约束最小化问题并不能保证找到严格分离器。感知机在间隔为零时仍更新，直至所有训练间隔严格为正。因此，\emph{零损失条件与算法终止条件并不相同}，线性可分条件也不会因为损失函数是凸函数而自动成立。

\section{感知机学习算法}
\label{sec:lh-perceptron-algorithm}

\subsection{原始形式与逐次更新}
\label{subsec:lh-perceptron-primal}

在没有样本跨越边界的局部区域内，负间隔样本集保持不变，此时损失的梯度为
\[
\nabla_w L=-\sum_{i:m_i<0}y_i x_i,
\qquad
\frac{\partial L}{\partial b}=-\sum_{i:m_i<0}y_i.
\]
若一次只选择一个间隔非正的样本 $(x_i,y_i)$，就得到原始形式（primal form）的更新
\begin{equation}
\label{eq:lh-perceptron-primal-update}
w\leftarrow w+\eta y_i x_i,
\qquad
b\leftarrow b+\eta y_i,
\qquad \eta>0.
\end{equation}
其中 $\eta$ 为固定学习率（learning rate）。教材采用不大于 $1$ 的步长，并在例题中取 $1$；本章的有限更新证明只要求它是固定正常数。

从优化角度看，这是逐样本的\emph{次梯度（subgradient）}更新。对凸函数 $h$，若向量 $g$ 满足对所有 $v$ 都有 $h(v)\geq h(u)+g^\top(v-u)$，则称 $g$ 为 $h$ 在 $u$ 处的次梯度，所有这样的向量组成次微分 $\partial h(u)$。令扩充参数 $\widetilde w=(w^\top,b)^\top$，扩充输入 $\widetilde x_i=(x_i^\top,1)^\top$，则
\[
\partial\ell_i(\widetilde w)=
\begin{cases}
\{-y_i\widetilde x_i\},&m_i<0,\\
\{-\lambda y_i\widetilde x_i:0\leq\lambda\leq1\},&m_i=0,\\
\{0\},&m_i>0.
\end{cases}
\]
边界处选择 $\lambda=1$ 就得到式~\eqref{eq:lh-perceptron-primal-update}，而选择 $\lambda=0$ 会允许算法停在零参数。教材用随机梯度下降（stochastic gradient descent）描述逐点修正过程；严格地说，在不可导位置应使用次梯度，且固定顺序遍历样本也能得到后面的有限更新保证，无须依赖随机抽样。

原始形式的输入是训练集 $T$、学习率 $\eta>0$ 和初值 $(w_0,b_0)$，输出是分类参数与式~\eqref{eq:lh-perceptron-model} 中的分类器。为使选点与停止条件明确，可采用下面的循环遍历实现。
\begin{enumerate}
  \item 初始化 $w=w_0$、$b=b_0$，常用 $w_0=0$、$b_0=0$。
  \item 开始一轮训练，将本轮更新次数置为零，依次访问 $i=1,\ldots,N$。
  \item 每访问一个样本，计算当前间隔 $m_i=y_i(w^\top x_i+b)$；若 $m_i\leq0$，执行式~\eqref{eq:lh-perceptron-primal-update} 并累计一次更新；否则保持参数不变。
  \item 完整一轮没有发生更新时停止；否则开始下一轮。在线性可分性未知时另设最大轮数，达到上限只能报告尚未找到严格分离器，不能判定数据必定不可分。
\end{enumerate}

一轮无更新意味着参数在整轮中固定，故所有样本在最终参数下的间隔均为正；仅仅在一轮中让每个点都曾经被正确分类，并不足以停止，因为后续更新可能破坏先前点的分类。教材允许直接选择任意间隔非正的点；循环遍历则保证只要还有这样的点，就会继续发现并处理它。

更新对当前点的作用可以精确写成
\begin{equation}
\label{eq:lh-perceptron-margin-increase}
y_i\bigl((w+\eta y_i x_i)^\top x_i+b+\eta y_i\bigr)
=m_i+\eta(\|x_i\|^2+1)>m_i.
\end{equation}
但一次更新未必把该点的间隔直接变为正。对其他样本 $(x_j,y_j)$，间隔改变量为 $\eta y_i y_j(x_i^\top x_j+1)$，可正可负，所以不能由式~\eqref{eq:lh-perceptron-margin-increase} 推出总损失每步下降，或所有样本同时得到改善。

\begin{example}
\label{ex:lh-perceptron-primal}
教材例 2.1：正例为 $x_1=(3,3)^\top$、$x_2=(4,3)^\top$，负例为 $x_3=(1,1)^\top$。从 $w_0=0$、$b_0=0$ 出发，取 $\eta=1$，用原始形式求一个感知机分类器。
\end{example}

\begin{solve}
每次从样本 $1$ 开始检查，选择第一个间隔非正的点更新，再重新检查。这样得到的更新点顺序为 $x_1,x_3,x_3,x_3,x_1,x_3,x_3$。下表中 $k$ 只计发生参数更新的次数。
\begin{center}
\begin{tabular}{cccc}
\toprule
$k$ & 更新所用样本 & $w_k^\top$ & $b_k$ \\
\midrule
0 & 初值 & $(0,0)$ & $0$ \\
1 & $x_1$ & $(3,3)$ & $1$ \\
2 & $x_3$ & $(2,2)$ & $0$ \\
3 & $x_3$ & $(1,1)$ & $-1$ \\
4 & $x_3$ & $(0,0)$ & $-2$ \\
5 & $x_1$ & $(3,3)$ & $-1$ \\
6 & $x_3$ & $(2,2)$ & $-2$ \\
7 & $x_3$ & $(1,1)$ & $-3$ \\
\bottomrule
\end{tabular}
\end{center}
例如第 $4$ 次更新前，$x_3$ 的间隔为 $-((1,1)(1,1)^\top-1)=-1$，故仍需修正；更新后 $w=0$、$b=-2$，两个正例又需要重新检查。这说明中间状态不必始终对应非退化的超平面，也不必保持先前的分类结果。

最终三个样本的间隔依次为 $3,4,1$，均严格为正，因此得到分离直线 $x^{(1)}+x^{(2)}-3=0$ 和分类器
\[
f(x)=\operatorname{sign}\bigl(x^{(1)}+x^{(2)}-3\bigr).
\]
这一解不是唯一的。例如 $2x^{(1)}+x^{(2)}-5=0$ 也能严格分离相同数据，其三个间隔分别为 $4,6,2$。感知机没有规定从多个分离器中选择最大间隔者，不同选点顺序或初值可能给出不同结果。
\end{solve}

\begin{notealert}[type=caution,title={损失值不能替代训练停止检查}]
例~\ref{ex:lh-perceptron-primal} 中初始总损失为 $0$，第一次更新后负例的得分为 $7$，总损失反而变为 $7$。这并不违反后面的收敛定理：该定理控制的是间隔非正所触发的更新次数，并不要求 $L$ 沿迭代单调递减。实现时应检查所有训练间隔是否严格为正，不能以“总损失已为零”替代算法的停止规则。
\end{notealert}

\subsection{算法的收敛性}
\label{subsec:lh-perceptron-convergence}

为统一处理权重与偏置，沿用扩充表示
\begin{equation}
\label{eq:lh-perceptron-augmentation}
\widetilde x_i=(x_i^\top,1)^\top\in\mathbb R^{d+1},
\qquad
\widetilde w=(w^\top,b)^\top\in\mathbb R^{d+1}.
\end{equation}
于是 $\widetilde w^\top\widetilde x_i=w^\top x_i+b$，更新成为 $\widetilde w\leftarrow\widetilde w+\eta y_i\widetilde x_i$。\emph[caution]{扩充参数的范数为 $\sqrt{\|w\|^2+b^2}$，与输入空间超平面距离公式中使用的 $\|w\|$ 不同。}

\begin{theorem}[感知机有限更新定理，Novikoff]
\label{thm:lh-perceptron-convergence}
设有限训练集 $T$ 线性可分，使用固定学习率 $\eta>0$，从 $\widetilde w_0=0$ 开始，只在 $y_i\widetilde w^\top\widetilde x_i\leq0$ 时更新。则存在单位向量 $\widetilde u\in\mathbb R^{d+1}$ 及常数 $\gamma>0$，使
\begin{equation}
\label{eq:lh-perceptron-separator-margin}
\|\widetilde u\|=1,
\qquad
y_i\widetilde u^\top\widetilde x_i\geq\gamma,
\quad i=1,\ldots,N.
\end{equation}
令 $R=\max_{1\leq i\leq N}\|\widetilde x_i\|$，则算法发生的更新总次数 $K$ 满足
\begin{equation}
\label{eq:lh-perceptron-mistake-bound}
K\leq\left(\frac{R}{\gamma}\right)^2.
\end{equation}
只要选点规则保证持续检查整个训练集，算法就会在有限次检查后得到所有训练间隔严格为正的参数。
\end{theorem}

\begin{proof}
先说明单位分离向量及正间隔的存在性。由线性可分性，存在非零扩充参数 $\widetilde v$，使 $y_i\widetilde v^\top\widetilde x_i>0$ 对所有 $i$ 成立。令 $\widetilde u=\widetilde v/\|\widetilde v\|$，并取
\[
\gamma=\min_{1\leq i\leq N}y_i\widetilde u^\top\widetilde x_i.
\]
因为样本数有限，有限个严格正数的最小值仍严格为正，故式~\eqref{eq:lh-perceptron-separator-margin} 成立。$\widetilde u$ 在这里仅是用于分析的一个单位分离向量，无须预先求得，更没有假设它使间隔达到最大。

用 $i_k$ 表示第 $k$ 次更新选中的样本编号。更新前有
\[
y_{i_k}\widetilde w_{k-1}^\top\widetilde x_{i_k}\leq0,
\qquad
\widetilde w_k=\widetilde w_{k-1}+\eta y_{i_k}\widetilde x_{i_k}.
\]
证明的关键是比较两个增长速度：参数在正确分离方向上的投影至少按 $k$ 线性增长，而参数范数至多按 $\sqrt{k}$ 增长。

对投影，由式~\eqref{eq:lh-perceptron-separator-margin} 得
\[
\widetilde w_k^\top\widetilde u
=\widetilde w_{k-1}^\top\widetilde u
+\eta y_{i_k}\widetilde u^\top\widetilde x_{i_k}
\geq\widetilde w_{k-1}^\top\widetilde u+\eta\gamma.
\]
从零初值递推可得
\begin{equation}
\label{eq:lh-perceptron-projection-bound}
\widetilde w_k^\top\widetilde u\geq k\eta\gamma.
\end{equation}
另一方面，由更新条件及 $\|\widetilde x_{i_k}\|\leq R$，
\[
\begin{aligned}
\|\widetilde w_k\|^2
&=\|\widetilde w_{k-1}\|^2
+2\eta y_{i_k}\widetilde w_{k-1}^\top\widetilde x_{i_k}
+\eta^2\|\widetilde x_{i_k}\|^2\\
&\leq\|\widetilde w_{k-1}\|^2+\eta^2R^2.
\end{aligned}
\]
继续递推得到
\begin{equation}
\label{eq:lh-perceptron-norm-bound}
\|\widetilde w_k\|^2\leq k\eta^2R^2.
\end{equation}
若 $k\geq1$，结合式~\eqref{eq:lh-perceptron-projection-bound}、式~\eqref{eq:lh-perceptron-norm-bound} 和 Cauchy--Schwarz 不等式，有
\[
k\eta\gamma
\leq\widetilde w_k^\top\widetilde u
\leq\|\widetilde w_k\|\|\widetilde u\|
\leq\eta R\sqrt{k}.
\]
除以正数 $\eta\sqrt{k}\gamma$ 并平方，得到 $k\leq(R/\gamma)^2$；$k=0$ 时同样满足该界。因此不可能出现无限次更新。若循环遍历的某一阶段仍有非正间隔样本，则至多再检查一整轮就会发现可更新样本；故更新次数有限必然导致某一整轮无更新，并在严格分离状态停止。
\end{proof}

式~\eqref{eq:lh-perceptron-mistake-bound} 中的 $K$ 是更新次数，包含落在边界时的更新，不是遍历轮数，也不是程序运行时间。证明不依赖误分点的具体选择顺序；对于给定的单位分离向量，$R/\gamma$ 越大，所得最坏情况上界越宽松。这里的 $\gamma$ 是扩充空间中相对于单位参数的间隔，不能直接与输入空间中除以 $\|w\|$ 的几何间隔混用。

上界不含固定学习率 $\eta$，这也可以从零初值的缩放性质理解：在相同选点规则下，将学习率由 $1$ 改为 $\eta>0$，每一步参数都恰为原参数的 $\eta$ 倍，间隔符号与更新点顺序不变。该结论针对本章没有额外阈值或正则项的更新规则，并不表示一般优化算法的行为都与学习率无关。

对任意固定非零初值，有限更新性也可由同一思路得到，但不能原样套用式~\eqref{eq:lh-perceptron-mistake-bound}。此时两条递推分别成为
\[
\widetilde w_k^\top\widetilde u
\geq\widetilde w_0^\top\widetilde u+k\eta\gamma,
\qquad
\|\widetilde w_k\|^2
\leq\|\widetilde w_0\|^2+k\eta^2R^2.
\]
若更新次数能够无限增长，左侧投影的线性下界最终会超过由范数给出的平方根上界，产生矛盾。因此仍会有限终止，只是次数界还依赖初值。

若训练集线性不可分，就不存在满足式~\eqref{eq:lh-perceptron-separator-margin} 的正数 $\gamma$，上述证明失去前提。持续遍历并坚持严格正间隔的停止条件时，算法无法终止；参数可能反复变化，不能把有限可分情形的结论推广过去。最大轮数只是计算预算，若要处理不可分数据，需要另行选择允许误差的目标或算法。即使训练集可分，所得分类器对未见数据的表现仍需独立评价。

\subsection{对偶形式与 Gram 矩阵}
\label{subsec:lh-perceptron-dual}

从 $w_0=0$、$b_0=0$ 出发，每次只增加 $\eta y_i x_i$ 和 $\eta y_i$。若第 $i$ 个样本共触发 $n_i$ 次更新，令 $\alpha_i=\eta n_i\geq0$，则任一迭代时刻都有
\begin{equation}
\label{eq:lh-perceptron-dual-representation}
w=\sum_{j=1}^{N}\alpha_jy_jx_j,
\qquad
b=\sum_{j=1}^{N}\alpha_jy_j.
\end{equation}
将其代入得分可得
\begin{equation}
\label{eq:lh-perceptron-dual-score}
s(x)=\sum_{j=1}^{N}\alpha_jy_jx_j^\top x+b.
\end{equation}
这就是\emph{对偶形式（dual form）}的基本想法：以样本线性组合的系数表示权重，用内积完成训练和预测。这里的“对偶”首先是一种等价参数表示，不需要另行声称已经推导出某个拉格朗日对偶优化问题。

\begin{definition}[Gram 矩阵]
\label{def:lh-perceptron-gram}
训练样本的\emph{Gram 矩阵（Gram matrix）}定义为
\begin{equation}
\label{eq:lh-perceptron-gram}
G=(G_{ij})_{N\times N},
\qquad G_{ij}=x_i^\top x_j.
\end{equation}
它是对称半正定矩阵。事实上，对任意 $c\in\mathbb R^N$，
\[
c^\top Gc
=\left\|\sum_{i=1}^{N}c_i x_i\right\|^2\geq0.
\]
若采用扩充输入，则相应矩阵为 $\widetilde G_{ij}=G_{ij}+1$。
\end{definition}

对偶算法输入训练集 $T$ 和固定学习率 $\eta>0$，输出系数 $\alpha$、偏置 $b$ 及式~\eqref{eq:lh-perceptron-dual-score} 对应的分类器。先计算 $G$，并初始化 $\alpha_1=\cdots=\alpha_N=0$、$b=0$；按原始形式相同的规则访问样本。当
\begin{equation}
\label{eq:lh-perceptron-dual-condition}
y_i\left(\sum_{j=1}^{N}\alpha_jy_jG_{ji}+b\right)\leq0
\end{equation}
时，执行
\begin{equation}
\label{eq:lh-perceptron-dual-update}
\alpha_i\leftarrow\alpha_i+\eta,
\qquad
b\leftarrow b+\eta y_i,
\end{equation}
其余系数不变。完整一轮没有更新时停止；可分性未知时使用与原始形式相同的轮数上限及未完成说明。由于 $G$ 对称，式~\eqref{eq:lh-perceptron-dual-condition} 中也可写 $G_{ij}$，但必须保留求和系数中的标签 $y_j$。

\begin{proposition}[两种形式的逐步等价性]
\label{prop:lh-perceptron-primal-dual}
若原始形式与对偶形式都从零初值开始，使用相同学习率，并按相同规则访问和选择样本，则每次更新后两者的 $w,b$ 都满足式~\eqref{eq:lh-perceptron-dual-representation}，得分及更新条件完全一致。
\end{proposition}

\begin{proof}
初始时所有参数为零，等式成立。假设某一步成立，则原始形式在样本 $x_i$ 上的得分为
\[
w^\top x_i+b
=\sum_{j=1}^{N}\alpha_jy_jG_{ji}+b,
\]
因而两种形式决定是否更新的条件相同。若更新，对偶形式只令 $\alpha_i$ 增加 $\eta$，故所表示的权重增加 $\eta y_i x_i$，偏置增加 $\eta y_i$，恰好与原始形式一致；不更新时等式保持。由归纳法，结论对所有迭代成立。
\end{proof}

偏置既可以作为独立变量递推，也可以每次由 $b=\sum_j\alpha_jy_j$ 恢复，两者不能重复累加。若直接使用扩充 Gram 矩阵，则得分应写成 $\sum_j\alpha_jy_j\widetilde G_{ji}$，此时已经包含偏置贡献，不能再加一个 $b$。零初始化是式~\eqref{eq:lh-perceptron-dual-representation} 的前提；一般初值时还需保留 $w_0,b_0$ 项。

\begin{example}
\label{ex:lh-perceptron-dual}
教材例 2.2：对例~\ref{ex:lh-perceptron-primal} 的三个样本，使用同一更新点顺序，以对偶形式求感知机模型。
\end{example}

\begin{solve}
直接计算内积得到
\[
G=
\begin{pmatrix}
18&21&6\\
21&25&7\\
6&7&2
\end{pmatrix}.
\]
取 $\eta=1$，按更新点顺序 $x_1,x_3,x_3,x_3,x_1,x_3,x_3$ 计算，有
\begin{center}
\begin{tabular}{ccccc}
\toprule
$k$ & $\alpha_1$ & $\alpha_2$ & $\alpha_3$ & $b$ \\
\midrule
0 & 0 & 0 & 0 & 0 \\
1 & 1 & 0 & 0 & 1 \\
2 & 1 & 0 & 1 & 0 \\
3 & 1 & 0 & 2 & $-1$ \\
4 & 1 & 0 & 3 & $-2$ \\
5 & 2 & 0 & 3 & $-1$ \\
6 & 2 & 0 & 4 & $-2$ \\
7 & 2 & 0 & 5 & $-3$ \\
\bottomrule
\end{tabular}
\end{center}
因此
\[
w=2x_1-5x_3=(1,1)^\top,
\qquad
b=2-5=-3,
\]
所得模型与例~\ref{ex:lh-perceptron-primal} 相同。$\alpha_2=0$ 表示该点在这次训练轨迹中没有触发更新，并不表示它的标签没有约束训练结果；算法停止时仍须核查它的间隔为正。
\end{solve}

对稠密 $d$ 维特征，原始形式一次得分计算和一次权重更新各需 $O(d)$ 运算。预计算 Gram 矩阵通常需 $O(N^2d)$ 运算、$O(N^2)$ 存储，之后按式~\eqref{eq:lh-perceptron-dual-condition} 计算一个训练点的得分需 $O(N)$ 运算，而系数更新只涉及一个分量。因此对偶形式并非无条件更快，是否合算取决于样本数、维数及存储条件。训练完成后，也可以先恢复 $w$ 再用原始形式预测；若保留系数表示，则预测只需对非零 $\alpha_j$ 对应的样本求和。

系数 $\alpha_i/\eta$ 记录样本在特定训练轨迹中的更新次数，反映其对这条轨迹的累积贡献，但不能单凭更新次数断定该点到最终超平面的距离更小。它还依赖样本访问顺序及中间参数。对偶形式揭示了内积的作用，为以后引入核函数提供入口；本章的等价性与收敛性论证仍针对这里明确给出的线性内积。

\section{章末习题与进一步理解}
\label{sec:lh-perceptron-exercises}

\subsection{异或问题的不可分性}
\label{subsec:lh-perceptron-xor}

\begin{example}
\label{ex:lh-perceptron-xor}
教材习题 2.1：将二元异或（exclusive OR，XOR）的输出 $1$ 编为标签 $+1$，输出 $0$ 编为标签 $-1$，得到
\[
(0,0)\mapsto-1,\qquad
(1,0)\mapsto+1,\qquad
(0,1)\mapsto+1,\qquad
(1,1)\mapsto-1.
\]
证明单个感知机不能表示这一分类规则。
\end{example}

\begin{solve}
若存在严格分离参数 $w=(w_1,w_2)^\top,b$，四个样本依次要求
\[
b<0,\qquad w_1+b>0,\qquad w_2+b>0,\qquad w_1+w_2+b<0.
\]
两个正例的约束相加得 $w_1+w_2+2b>0$，两个负例的约束相加则得 $w_1+w_2+2b<0$，矛盾。因此数据不满足定义~\ref{def:lh-linear-separability}。

即使允许正例落在边界，并按定义~\ref{def:lh-perceptron} 将零得分判为 $+1$，也不能表示 XOR：此时两个正例仍要求 $w_1+w_2+2b\geq0$，两个负例仍要求该量严格小于零，矛盾不变。这个限制来自原始特征空间中单个线性边界的表达能力，增加训练轮数无法消除它。
\end{solve}

\subsection{自行构造训练例子}
\label{subsec:lh-perceptron-constructed-example}

\begin{example}
\label{ex:lh-perceptron-constructed}
教材习题 2.2 要求构造从训练集求感知机模型的例子。取一维数据 $x_1=2,y_1=+1$，$x_2=1,y_2=-1$，从 $w_0=b_0=0$ 开始，令 $\eta=1$，每次重新按 $x_1,x_2$ 的顺序寻找间隔非正的样本。求迭代结果，并核查是否严格分离。
\end{example}

\begin{solve}
一维时参数 $w$ 为实数。逐步计算得到
\begin{center}
\begin{tabular}{cccc}
\toprule
$k$ & 更新所用样本 & $w_k$ & $b_k$ \\
\midrule
0 & 初值 & 0 & 0 \\
1 & $x_1$ & 2 & 1 \\
2 & $x_2$ & 1 & 0 \\
3 & $x_2$ & 0 & $-1$ \\
4 & $x_1$ & 2 & 0 \\
5 & $x_2$ & 1 & $-1$ \\
6 & $x_2$ & 0 & $-2$ \\
7 & $x_1$ & 2 & $-1$ \\
8 & $x_2$ & 1 & $-2$ \\
9 & $x_1$ & 3 & $-1$ \\
10 & $x_2$ & 2 & $-2$ \\
11 & $x_2$ & 1 & $-3$ \\
12 & $x_1$ & 3 & $-2$ \\
13 & $x_2$ & 2 & $-3$ \\
\bottomrule
\end{tabular}
\end{center}
例如 $(w,b)=(1,-2)$ 时，正例的得分为零，必须继续更新；$(w,b)=(2,-2)$ 时，负例的得分为零，也必须更新。最终得到 $f(x)=\operatorname{sign}(2x-3)$，两个训练间隔均为 $1$，严格分离点为 $x=3/2$。这一例子突出边界处理：若将判据误写为 $m_i<0$，从零初值出发就不会发生任何更新。
\end{solve}

\subsection{线性可分与两类凸包不相交}
\label{subsec:lh-perceptron-convex-hulls}

\begin{definition}[有限点集的凸包]
\label{def:lh-perceptron-convex-hull}
对非空有限点集 $A=\{a_1,\ldots,a_r\}\subset\mathbb R^d$，其\emph{凸包（convex hull）}为所有凸组合构成的集合：
\[
\operatorname{conv}(A)=
\left\{\sum_{j=1}^{r}\lambda_j a_j:
\lambda_j\geq0,\ \sum_{j=1}^{r}\lambda_j=1\right\}.
\]
它是非空紧凸集：非负且和为 $1$ 的系数集合是有限维闭有界单纯形，凸包是该紧集在线性映射下的像；两组凸组合再作凸组合仍为凸组合。
\end{definition}

\begin{theorem}[两类凸包判据]
\label{thm:lh-perceptron-convex-hulls}
教材习题 2.3：设有限训练集的正例点集 $A=\{x_i:y_i=+1\}$ 与负例点集 $B=\{x_i:y_i=-1\}$ 均非空。则训练集线性可分，当且仅当
\begin{equation}
\label{eq:lh-perceptron-disjoint-hulls}
\operatorname{conv}(A)\cap\operatorname{conv}(B)=\varnothing.
\end{equation}
\end{theorem}

\begin{proof}
记 $C_+=\operatorname{conv}(A)$、$C_-=\operatorname{conv}(B)$。先证必要性。若 $s(x)=w^\top x+b$ 严格分离两个点集，则任意正例凸组合 $a=\sum_j\lambda_j a_j$ 满足
\[
s(a)=\sum_j\lambda_j s(a_j)>0,
\]
这里用了系数非负、和为 $1$，且至少有一个系数为正。同理，每个 $c\in C_-$ 都满足 $s(c)<0$。一个点不可能同时具有正得分和负得分，故两个凸包不相交。

再证充分性。假设两个凸包不相交。由于它们是非空紧集，连续函数 $(a,c)\mapsto\|a-c\|$ 在 $C_+\times C_-$ 上取得最小值。取最近点对 $p\in C_+$、$q\in C_-$，则 $p\neq q$，令
\[
w=p-q\neq0,
\qquad
b=-\frac12 w^\top(p+q).
\]
下面证明垂直于 $p-q$ 且经过两点中点的超平面严格分离两个凸包。

任取 $a\in C_+$，凸性保证 $p+t(a-p)\in C_+$ 对 $0\leq t\leq1$ 成立。由最近点性质，
\[
\|p+t(a-p)-q\|^2\geq\|p-q\|^2.
\]
展开后对 $t>0$ 除以 $t$，再令 $t\downarrow0$，得到 $w^\top(a-p)\geq0$。同样，对任意 $c\in C_-$，沿 $q+t(c-q)$ 变化有
\[
\|p-q-t(c-q)\|^2\geq\|p-q\|^2,
\]
从而 $w^\top(c-q)\leq0$。因此
\[
\begin{aligned}
w^\top a+b
&\geq w^\top p-\tfrac12w^\top(p+q)
=\tfrac12\|p-q\|^2>0,\\
w^\top c+b
&\leq w^\top q-\tfrac12w^\top(p+q)
=-\tfrac12\|p-q\|^2<0.
\end{aligned}
\]
两个凸包中的全部点均被严格分开，原训练样本作为凸包中的点也被严格分开，故训练集线性可分。
\end{proof}

有限性在充分性证明中保证了凸包紧性及最近点对的存在；不能不加条件地把这一论证照搬到任意不相交集合。对 XOR，两个正例的凸包与两个负例的凸包都包含 $(1/2,1/2)^\top$，由定理~\ref{thm:lh-perceptron-convex-hulls} 可再次看出不可分。若同一个特征向量在训练集中具有相反标签，它本身就同时属于两个凸包，也立即排除了严格线性分离。

本章的核心联系是：模型用仿射得分的符号分类，学习规则以非正函数间隔触发更新，有限可分性通过单位分离方向上的正间隔给出更新次数界，对偶形式则将相同轨迹表示为样本内积及非负更新系数。损失的连续凸性、算法有限终止和未知样本上的泛化是三个不同层面的性质，使用时应分别核对其定义与条件。

% 教材来源：机器学习/textbooks/机器学习（分章原文）/机器学习-3.pdf，第 49--58 页。





